% Part: first-order-logic
% Chapter: introduction
% Section: first-order-logic

\documentclass[../../../include/open-logic-section]{subfiles}

\begin{document}

\olfileid{fol}{int}{fol}

\olsection{منطق الرتبة الأولى}

لعل منطق الرتبة الأولى مألوف لديك من مدخلك الأول إلى المنطق
الصوري.\footnote{بل إننا نفترض، إلى حد بعيد، أنه مألوف لديك!{} فإن لم
  يكن كذلك، أمكنك مراجعة كتاب أكثر أولية، مثل \emph{forall x}
  \citep{Magnus2021}.} وربما تعرفه باسم «منطق الكم» أو «منطق
المحمولات». ومنطق الرتبة الأولى، قبل كل شيء، لغة صورية. وهذا يعني أن له
معجمًا معينًا، وأن تعبيراته سلاسل مؤلفة من عناصر هذا المعجم. لكن ليست
كل سلسلة مسموحًا بها. فللتعبيرات المسموح بها أنواع مختلفة: الحدود،
وال!!{formula}s، وال!!{sentence}s. ويتركز اهتمامنا أساسًا على
ال!!{sentence}s في منطق الرتبة الأولى؛ فهي تزودنا بنظير صوري لجمل اللغة
الطبيعية، ويمكننا أن نسأل عنها الأسئلة التي تهم المناطقة عادة. ومن ذلك:
\begin{itemize}
    \item هل تلزم $!B$ منطقيًا عن~$!A$؟
    \item هل $!A$ صادقة منطقيًا، أم كاذبة منطقيًا، أم عارضة؟
    \item هل $!A$ و$!B$ متكافئتان؟
\end{itemize}

هذه أسئلة تتعلق في المقام الأول بـ«معنى» ال!!{sentence}s في منطق الرتبة
الأولى. فمثلًا، يحلل الفيلسوف السؤال عما إذا كانت $!B$ تلزم منطقيًا
عن~$!A$ على النحو الآتي: هل توجد حالة تصدق فيها $!A$ وتكذب~$!B$ (فلا
تلزم $!B$ عن~$!A$)، أم أن كل حالة تجعل $!A$ صادقة تجعل~$!B$ صادقة
أيضًا (فتلزم $!B$ عن~$!A$)؟ لكننا لم نحدد بعد ما «الحالة»؛ فهذه مهمة
\emph{الدلالة}. وتقدم دلالة منطق الرتبة الأولى نموذجًا دقيقًا رياضيًا
لفكرة الفيلسوف الحدسية عن «الحالة»، وكذلك---وهذا مهم---لما يعنيه أن
تكون !!a{sentence}~$!A$ \emph{صادقة في} حالة. ونسمي النموذج الرياضي
الدقيق الذي سنطوره !!a{structure}. أما العلاقة التي تضبط معنى «صادقة
في» فتسمى علاقة \emph{الإرضاء}. ومن ثم سنعرّف معنى «تُرضى $!A$
في~$\Struct{M}$» (ورمز ذلك: $\Sat{M}{!A}$) لل!!{sentence}s~$!A$
وال!!{structure}s~$\Struct{M}$. وبعد ذلك يمكننا أيضًا أن نعطي تعريفات
دقيقة للمصطلحات الدلالية الأخرى، مثل «تلزم عن» و«صادقة منطقيًا». وستتيح
هذه التعريفات أن نحسم، بالدقة الرياضية مرة أخرى، ما إذا كان، مثلًا،
$\lforall[x][(!A(x) \lif !B(x))],
\lexists[x][!A(x)] \Entails \lexists[x][!B(x)]$. والجواب طبعًا «نعم».
فإذا كنت قد تدربت من قبل على ترميز جمل اللغة الطبيعية بمنطق الرتبة
الأولى، فستتعرف في ذلك، مثلًا، ترميز الحجة: «كل النمل حشرات، ويوجد نمل،
إذن توجد حشرات». وهذه حجة صالحة بوضوح، ولذلك ينبغي أن يعطي نموذجنا
الرياضي لعبارة «تلزم عن» في لغتنا الصورية الجواب نفسه.

ومن الموضوعات الأخرى التي لعلها مألوفة لديك من مدخلك الأول إلى المنطق
الصوري وجود \emph{!!{derivation}s}. فإذا درست مقررًا أول في المنطق
الصوري، فلا بد أن مدرسك قد دربك على إيجاد مثل هذه ال!!{derivation}s،
وربما حتى على إيجاد !!a{derivation} يبيّن صحة الاستلزام السابق. وتوجد
طرائق كثيرة مختلفة لعرض ال!!{derivation}s: فربما استعملت ما يسمى
«الاستنباط الطبيعي» أو «أشجار الصدق»، وثمة طرائق أخرى كثيرة. والغرض من
أنساق ال!!{derivation} أن توفر أدوات يمكن بها الإجابة عن أسئلة المناطقة
السابقة. فمثلًا، !!a{derivation} في الاستنباط الطبيعي يكون فيه
$\lforall[x][(!A(x) \lif !B(x))]$ و$\lexists[x][!A(x)]$ مقدمتين،
و$\lexists[x][!B(x)]$ نتيجةً (السطر الأخير)، \emph{يثبت} أن
$\lexists[x][!B(x)]$ تلزم منطقيًا عن
$\lforall[x][(!A(x) \lif !B(x))]$ و$\lexists[x][!A(x)]$.

ولكن لماذا يصح ذلك؟ لا يبدو لأنساق ال!!{derivation}، في ظاهر الأمر،
أي صلة بالدلالة: فإعطاء !!a{derivation} صوري لا يعدو ترتيب الرموز على
وجوه معينة تحكمها قواعد؛ ولا يرد فيها ذكر «الحالات» أو «الصدق في»
أصلًا. ولا بد من إقامة الصلة بين أنساق ال!!{derivation} والدلالة ببحث
فوق منطقي. والمطلوب هو برهان رياضي، مثلًا، على أن وجود
!!a{derivation} صوري لـ~$\lexists[x][!B(x)]$ من المقدمتين
$\lforall[x][(!A(x) \lif !B(x))]$ و$\lexists[x][!A(x)]$ ممكن إذا
وفقط إذا كانت $\lforall[x][(!A(x) \lif !B(x))]$ و
$\lexists[x][!A(x)]$ تستلزمان معًا~$\lexists[x][!B(x)]$. لكن قبل أن
يتسنى ذلك، يلزم عمل دقيق شاق لضبط تعريفي النحو والدلالة.

\end{document}
